\section{Download the project}
\label{sec:clone}

\subsection{Background}

Copying a repository from GitHub to your own computer is called \term{cloning}. The clone is an
ordinary folder that contains every file of the project together with its history. This guide
places it in the folder \file{C:\bs ee105\bs camera-to-latex}, and every later command assumes
that location.

\rationale{The folder is deliberately short, free of spaces, and outside the \file{Desktop} and
\file{Documents} folders. On many laptops those two folders are synchronized by OneDrive, which
tries to upload each of the tens of thousands of small files that the Python environment
contains; an installation inside them can take ten times as long. Short paths also keep some
very long file names of TensorFlow within the length that Windows permits.}

\subsection{Clone the repository}

\step{Create the folder for course work}
In a PowerShell window run:
\begin{ps}
mkdir C:\ee105
\end{ps}
\begin{out}
    Directory: C:\

Mode                 LastWriteTime         Length Name
----                 -------------         ------ ----
d-----        11/24/2026   9:52 AM                ee105
\end{out}
\checkpoint{The listing shows the new folder \file{ee105}. If PowerShell says that the item
already exists, the folder was created earlier, which is equally good.}

\step{Go into the folder}
\begin{ps}
cd C:\ee105
\end{ps}
\checkpoint{The prompt now reads \texttt{PS C:\bs ee105>}.}

\step{Clone}\label{st:gitclone}
\begin{ps}
git clone https://github.com/ashahvos/camera-to-latex.git
\end{ps}
If the repository is private, a window titled \ui{Connect to GitHub} opens within a few
seconds. Click \ui{Sign in with your browser}. The browser opens a GitHub page; sign in with
the account of Section~\ref{sec:github} if asked, click \ui{Authorize} (the page names the
\emph{Git Credential Manager}), and return to PowerShell when the page says that
authentication succeeded. Windows remembers this sign-in, so it happens only once.
\begin{out}
Cloning into 'camera-to-latex'...
remote: Enumerating objects: 82, done.
remote: Counting objects: 100% (82/82), done.
remote: Compressing objects: 100% (74/74), done.
Receiving objects: 100% (82/82), 702.11 KiB | 6.21 MiB/s, done.
Resolving deltas: 100% (5/5), done.
\end{out}
\checkpoint{The output ends with lines that say \emph{done}, and no line begins with
\emph{fatal}. The numbers will differ.}

\begin{iffail}[If the output says ``Repository not found'']
GitHub hides a private repository from accounts that have no access, so this message means one
of three things: the invitation of Section~\ref{sec:github} has not been accepted; you signed in
to a different GitHub account than the invited one; or the address was mistyped. Open
\repourl{} in the browser. If the page loads while you
are signed in, the first two causes are excluded, and you should copy the command again from
this PDF.
\end{iffail}

\step{Enter the project folder and look at it}
\begin{ps}
cd C:\ee105\camera-to-latex
\end{ps}
\begin{ps}
Get-ChildItem
\end{ps}
\checkpoint{The prompt shows the folder \file{C:\bs ee105\bs camera-to-latex}, and the listing
contains the names of Table~\ref{tab:repo}, among them \file{config.py}, \file{requirements.txt},
\file{pipeline} and \file{tools}.}

The same folder can be opened in File Explorer at any time with the command
\code|explorer .| (the word \emph{explorer}, a space, and a full stop, which stands for ``the
current folder'').

\subsection{Introduce yourself to Git}

Git labels every commit with the name and email address of its author, and refuses to create a
commit until it knows them. The four commands of this subsection are given once per computer.

\step{Set your name}
In the next command, replace the yellow placeholder by your first and last name. Keep the
quotation marks.
\begin{center}
\texttt{git config -{}-global user.name "}\ph{YOUR-NAME}\texttt{"}
\end{center}
For a team member named Tommy Trojan the finished command would be
\begin{out}[Example only; type your own name]
git config --global user.name "Tommy Trojan"
\end{out}

\step{Set your email address}
Replace the placeholder by the email address of your GitHub account, again inside the quotation
marks.
\begin{center}
\texttt{git config -{}-global user.email "}\ph{YOUR-EMAIL}\texttt{"}
\end{center}
\begin{out}[Example only; type your own address]
git config --global user.email "ttrojan@usc.edu"
\end{out}

\step{Choose how Git combines and describes changes}
The next two commands are typed exactly as printed. The first tells Git to combine your work
with the team's by a \term{merge}, the simplest method. The second makes Git open Notepad,
rather than a terminal text editor that is difficult to leave, on the rare occasions when it
needs a sentence from you.
\begin{ps}
git config --global pull.rebase false
\end{ps}
\begin{ps}
git config --global core.editor notepad
\end{ps}

\step{Check the four settings}
\begin{ps}
git config --global --list
\end{ps}
\begin{out}[What you should see (other lines may also appear)]
user.name=Tommy Trojan
user.email=ttrojan@usc.edu
pull.rebase=false
core.editor=notepad
\end{out}
\checkpoint{The four lines are present, with your own name and address in the first two.}
